\documentclass[12pt,a4paper]{article}
\usepackage[utf8]{inputenc}
\usepackage[indonesian]{babel}
\usepackage{amsmath}
\usepackage{amsfonts}
\usepackage{amssymb}
\usepackage{geometry}

\geometry{margin=2.5cm}

\title{Menghitung Volume Benda Putar}
\author{Ahmad Fakhrul Bawani}
\date{}

\begin{document}

\maketitle

\section*{Soal}

Find the volume of the solid generated by revolving the region bounded by $y = 4\sqrt{x}$ and the lines $y = 8$ and $x = 0$ about:
\begin{center}
  \textbf{a.} the $x$-axis. \qquad \textbf{b.} the $y$-axis. \qquad \textbf{c.} the line $y = 8$. \qquad \textbf{d.} the line $x = 4$.
\end{center}

---

\subsection*{a. Diputar terhadap Sumbu-$x$ (Poros Horizontal)}
Menggunakan \textbf{Metode Cincin (Washer Method)} terhadap variabel $x$ dari $x = 0$ hingga $x = 4$:
\begin{itemize}
  \item Jari-jari luar (Outer radius): $R(x) = 8$
  \item Jari-jari dalam (Inner radius): $u(x) = 4\sqrt{x}$
  \item \textbf{Setup Integral:}
    \begin{align*}
      V &= \pi \int_{0}^{4} \left[ [R(x)]^2 - [u(x)]^2 \right] \, dx \\
      V &= \pi \int_{0}^{4} \left[ 8^2 - (4\sqrt{x})^2 \right] \, dx = \pi \int_{0}^{4} (64 - 16x) \, dx
    \end{align*}
  \item \textbf{Perhitungan:}
    \begin{equation*}
      V = \pi \left[ 64x - 8x^2 \right]_{0}^{4} = \pi \left[ 64(4) - 8(4)^2 \right] = \pi (256 - 128) = 128\pi
    \end{equation*}
\end{itemize}

\subsection*{b. Diputar terhadap Sumbu-$y$ (Poros Vertikal)}
Menggunakan \textbf{Metode Cakram (Disk Method)} terhadap variabel $y$ dari $y = 0$ hingga $y = 8$:
\begin{itemize}
  \item Jari-jari cakram horizontal: $R(y) = x = \frac{y^2}{16}$
  \item \textbf{Setup Integral:}
    \begin{equation*}
      V = \pi \int_{0}^{8} [R(y)]^2 \, dy = \pi \int_{0}^{8} \left(\frac{y^2}{16}\right)^2 \, dy = \pi \int_{0}^{8} \frac{y^4}{256} \, dy
    \end{equation*}
  \item \textbf{Perhitungan:}
    \begin{equation*}
      V = \frac{\pi}{256} \left[ \frac{1}{5}y^5 \right]_{0}^{8} = \frac{\pi}{256} \cdot \frac{32768}{5} = \frac{128\pi}{5}
    \end{equation*}
\end{itemize}

\subsection*{c. Diputar terhadap Garis $y = 8$ (Poros Horizontal)}
Menggunakan \textbf{Metode Cakram (Disk Method)} terhadap variabel $x$ dari $x = 0$ hingga $x = 4$:
\begin{itemize}
  \item Jari-jari cakram vertikal: $R(x) = 8 - 4\sqrt{x}$
  \item \textbf{Setup Integral:}
    \begin{align*}
      V &= \pi \int_{0}^{4} [R(x)]^2 \, dx = \pi \int_{0}^{4} (8 - 4\sqrt{x})^2 \, dx \\
      V &= \pi \int_{0}^{4} (64 - 64x^{\frac{1}{2}} + 16x) \, dx
    \end{align*}
  \item \textbf{Perhitungan:}
    \begin{align*}
      V &= \pi \left[ 64x - \frac{128}{3}x^{\frac{3}{2}} + 8x^2 \right]_{0}^{4} \\
      &= \pi \left( 64(4) - \frac{128}{3}(8) + 8(16) \right) \\
      &= \pi \left( 256 - \frac{1024}{3} + 128 \right) = \pi \left( 384 - \frac{1024}{3} \right) = \frac{128\pi}{3}
    \end{align*}
\end{itemize}

\subsection*{d. Diputar terhadap Garis $x = 4$ (Poros Vertikal)}
Menggunakan \textbf{Metode Cincin (Washer Method)} terhadap variabel $y$ dari $y = 0$ hingga $y = 8$:
\begin{itemize}
  \item Jari-jari luar: $R(y) = 4 - 0 = 4$
  \item Jari-jari dalam: $u(y) = 4 - \frac{y^2}{16}$
  \item \textbf{Setup Integral:}
    \begin{align*}
      V &= \pi \int_{0}^{8} \left[ [R(y)]^2 - [u(y)]^2 \right] \, dy \\
      V &= \pi \int_{0}^{8} \left[ 4^2 - \left(4 - \frac{y^2}{16}\right)^2 \right] \, dy \\
      V &= \pi \int_{0}^{8} \left[ 16 - \left(16 - \frac{y^2}{2} + \frac{y^4}{256}\right) \right] \, dy = \pi \int_{0}^{8} \left( \frac{y^2}{2} - \frac{y^4}{256} \right) \, dy
    \end{align*}
  \item \textbf{Perhitungan:}
    \begin{align*}
      V &= \pi \left[ \frac{1}{6}y^3 - \frac{1}{1280}y^5 \right]_{0}^{8} \\
      &= \pi \left( \frac{512}{6} - \frac{32768}{1280} \right) = \pi \left( \frac{256}{3} - \frac{128}{5} \right) \\
      &= \pi \left( \frac{1280 - 384}{15} \right) = \frac{896\pi}{15}
    \end{align*}
\end{itemize}

\subsection*{Kesimpulan Hasil Jawaban}
\begin{itemize}
  \item \textbf{a.} Volume terhadap sumbu-$x$ adalah \textbf{$128\pi$} cubic units.
  \item \textbf{b.} Volume terhadap sumbu-$y$ adalah \textbf{$\dfrac{128\pi}{5}$} cubic units.
  \item \textbf{c.} Volume terhadap garis $y = 8$ adalah \textbf{$\dfrac{128\pi}{3}$} cubic units.
  \item \textbf{d.} Volume terhadap garis $x = 4$ adalah \textbf{$\dfrac{896\pi}{15}$} cubic units.
\end{itemize}

\end{document}